\section{Cross-Fitted Orthogonal Estimation}

Partition the auction sample into folds $I_1,\ldots,I_K$. For each fold $k$,
train the WGAN-GP and estimate all nuisance functions using observations outside
$I_k$. Denote the resulting nuisance estimate by $\widehat\eta^{(-k)}$.

Sample splitting by itself is not sufficient for DML in this application. The
fold-averaged plug-in $\overline r_{\mathrm{FA}}$ introduced in Section~5 uses
the same training partitions but does not evaluate an observation-level score
on $I_k$. It is therefore an ensemble benchmark. Cross-fitting acquires its
inferential role only when each held-out observation enters the score below.

Let $\Phi$ and $\phi$ denote the standard normal CDF and density. For fixed
bandwidth $h$, define the observable signals
\begin{equation}
S_h(B;\theta)=\Phi\!\left(\frac{\theta-B}{h}\right),
\qquad
K_h(B;\theta)=\frac{1}{h}\phi\!\left(\frac{\theta-B}{h}\right),
\end{equation}
and the local nuisance vector
$\eta_h(\theta)=(A_h(\theta),G_h(\theta))$, where
$A_h(\theta)=\mathbb E[S_h(B;\theta)]$ and
$G_h(\theta)=\mathbb E[K_h(B;\theta)]$. The regularized reserve moment is
\begin{equation}
m_h(\theta,A,G)=1-A^{1/N}-\frac{\theta}{N}A^{1/N-1}G.
\label{eq:regularized-reserve-moment}
\end{equation}
Fix a positive compact search interval
$\Theta=[\underline\theta,\overline\theta]\subset(0,\infty)$ before estimation.
Define $R_h(\theta)=\theta\{1-A_h(\theta)^{1/N}\}$, so that
$R_h'(\theta)=m_h(\theta,A_h(\theta),G_h(\theta))$. Under the single-crossing
condition stated below, the fixed-bandwidth target $\theta_{0,h}$ is its unique
interior maximizer on $\Theta$. It need not equal the exact optimal reserve
$r_0$: Gaussian smoothing is applied to the maximum distribution before
inversion, not to the original valuation distribution. Approximation of $r_0$
requires a separate shrinking-bandwidth argument.

The derivatives needed for the locally robust correction are
\begin{align}
m_A(\theta,A,G)
&=-\frac{1}{N}A^{1/N-1}
-\frac{\theta}{N}\left(\frac{1}{N}-1\right)A^{1/N-2}G,\\
m_G(\theta,A,G)&=-\frac{\theta}{N}A^{1/N-1}.
\end{align}
The observation-level score is
\begin{equation}
\begin{split}
\psi_h(B;\theta,A,G)
={}&m_h(\theta,A,G)
+m_A(\theta,A,G)\{S_h(B;\theta)-A\}\\
&+m_G(\theta,A,G)\{K_h(B;\theta)-G\}.
\end{split}
\label{eq:orthogonal-reserve-score}
\end{equation}
Because the two residuals have mean zero at the truth, differentiation of the
score expectation adds $-m_A$ and $-m_G$ to the derivatives of the first term.
Hence
\begin{equation}
\left.\partial_{(A,G)}
\mathbb E\left[\psi_h(B;\theta_{0,h},A,G)\right]
\right|_{(A,G)=\eta_{0,h}(\theta_{0,h})}=0.
\end{equation}

\begin{proposition}[Orthogonality and remainder]
\label{prop:orthogonal-remainder}
At any fixed $\theta$, let $\eta_0=(A_h(\theta),G_h(\theta))$. Then
\begin{equation}
\mathbb E_0\psi_h(B;\theta,\eta)
=m_h(\theta,\eta)+\nabla_\eta m_h(\theta,\eta)^{\mathsf T}(\eta_0-\eta).
\end{equation}
On a compact nuisance region with $A$ bounded away from zero, this implies
\begin{equation}
\left|\mathbb E_0\psi_h(B;\theta,\eta)-m_h(\theta,\eta_0)\right|
\leq C_h\{(A-A_h)^2+|A-A_h|\,|G-G_h|\}.
\label{eq:anisotropic-remainder}
\end{equation}
If $A=A_h$, the expectation is correct for any finite $G$. Correct $G$ alone
does not generally have this property. The score is therefore Neyman
orthogonal and locally robust, but not a two-sided doubly robust AIPW score.
\end{proposition}
For uniform nuisance errors $\delta_{A,n}$ and $\delta_{G,n}$, the population
remainder is therefore negligible at the root-$n$ scale if
$\delta_{A,n}^2+\delta_{A,n}\delta_{G,n}=o_p(n^{-1/2})$.
This condition must be accompanied by score consistency and the usual
stochastic and root controls; it is not a license to substitute arbitrary
quantities for either local moment.

The estimator solves
\begin{equation}
\frac{1}{n}\sum_{k=1}^K\sum_{i\in I_k}
\psi_h\!\left(B_i;\widehat\theta_h,
\widehat A_h^{(-k)}(\widehat\theta_h),
\widehat G_h^{(-k)}(\widehat\theta_h)\right)=0.
\end{equation}

Figure~\ref{fig:estimation-workflow} summarizes the calibrated implementation.
The training observations supply both the WGAN and empirical local moments;
the held-out observations enter only through the score signals. Scores are
pooled across folds before solving for a common reserve estimate.

\begin{figure}[tbp]
\centering
\begin{tikzpicture}[
  box/.style={draw=black!55, rounded corners=3pt, line width=0.5pt,
    align=center, text width=5.2cm, minimum height=1.05cm,
    inner sep=5pt, font=\linespread{1}\fontsize{10}{12}\selectfont},
  flow/.style={-{Stealth[length=2mm]}, draw=black!65, line width=0.6pt}
]
\node[box] (train) at (0,0)
  {Training-fold maxima\\$B_j,\ j\in I_k^c$};
\node[box] (gan) at (-3.4,-1.8)
  {Train WGAN-GP\\Generate synthetic maxima};
\node[box] (empirical) at (3.4,-1.8)
  {Empirical local moments\\$\widehat A_{E,h}^{(-k)},\ \widehat G_{E,h}^{(-k)}$};
\node[box] (ganmoments) at (-3.4,-3.6)
  {WGAN local moments\\$\widehat A_{W,h}^{(-k)},\ \widehat G_{W,h}^{(-k)}$};
\node[box] (heldout) at (3.4,-3.6)
  {Held-out maxima\\$B_i,\ i\in I_k$};
\node[box] (calibration) at (-3.4,-5.4)
  {Calibrate local moments\\$\widehat A_h^{(-k)},\ \widehat G_h^{(-k)}$};
\node[box] (signals) at (3.4,-5.4)
  {Observed score signals\\$S_h(B_i;\theta),\ K_h(B_i;\theta)$};
\node[box] (score) at (0,-7.2)
  {Evaluate orthogonal scores\\on held-out observations};
\node[box] (pool) at (0,-9.0)
  {Rotate over three folds\\Pool held-out scores};
\node[box] (root) at (0,-10.8)
  {Unique downward-root diagnostic\\Estimate $\widehat\theta_h$ or report failure};
\draw[flow] (train.south) -- ++(0,-0.2) -| (gan.north);
\draw[flow] (train.south) -- ++(0,-0.2) -| (empirical.north);
\draw[flow] (gan) -- (ganmoments);
\draw[flow] (ganmoments) -- (calibration);
\draw[flow] (empirical.west) -- (0,-1.8) -- (0,-4.5)
  -- (-3.4,-4.5) -- (calibration.north);
\draw[flow] (heldout) -- (signals);
\draw[flow] (calibration.south) -- ++(0,-0.2)
  -| ([xshift=-10mm]score.north);
\draw[flow] (signals.south) -- ++(0,-0.2)
  -| ([xshift=10mm]score.north);
\draw[flow] (score) -- (pool);
\draw[flow] (pool) -- (root);
\end{tikzpicture}
\caption{Cross-fitted estimation with calibrated generative local moments.
Synthetic maxima are formed from $N$ independently generated bidder values.
For each candidate reserve, the training-fold WGAN and empirical moments are
combined using equation~\eqref{eq:local-calibration}. Training and tuning use
only $I_k^c$. The three-fold implementation pools observation-level
scores, rather than fold-specific reserve estimates. The final estimate targets
the fixed-bandwidth reserve $\theta_{0,h}$ and is reported only when the
prespecified root, overlap, and slope checks pass.}
\label{fig:estimation-workflow}
\end{figure}

Write $\widehat\Psi_n(\theta)$ for the pooled sample score in the preceding
equation. For the population moment and its derivative, use
\[
m_{0,h}(\theta)=m_h(\theta,A_h(\theta),G_h(\theta)),
\qquad J_h=m_{0,h}'(\theta_{0,h}).
\]
The formal estimator retains an interior downward
crossing only when it is the unique root on $\Theta$, the endpoint scores have
the prescribed positive and negative signs, and its total score derivative
is negative and separated from zero. Bracket refinement solves the score to
residual $o(n^{-1/2})$. Missing roots, multiple roots, boundary roots, active
CDF clipping, or a nearly zero derivative are reported as inference failures,
not replaced by a minimum-absolute-score point with a Wald interval.
The implementation scans a prespecified grid and refines detected crossings;
its finite-sample root count is a diagnostic, not a certificate excluding
every possible hidden crossing between grid points. The theorem below uses
uniform derivative convergence and a single population root to justify this
rule asymptotically, without requiring such a finite-sample certificate.

Let $I_k^c$ denote the training observations for fold $k$ and
$n_{-k}=|I_k^c|$. The empirical-local learner is
\begin{equation}
\widehat A_h^{(-k)}(\theta)
=\frac{1}{n_{-k}}\sum_{j\in I_k^c}S_h(B_j;\theta),
\qquad
\widehat G_h^{(-k)}(\theta)
=\frac{1}{n_{-k}}\sum_{j\in I_k^c}K_h(B_j;\theta).
\label{eq:empirical-local-nuisance}
\end{equation}

\begin{assumption}[Fixed-bandwidth regularity]
\label{ass:fixed-h}
The sampling restrictions in Section~3 hold, and $h>0$, $N$, and $K\geq2$
are fixed. The following conditions apply.
\begin{enumerate}
\item The interval $\Theta$ is fixed and positive. The continuous function
$m_{0,h}$ has exactly one zero $\theta_{0,h}$ in its interior,
$m_{0,h}(\underline\theta)>0$, and $m_{0,h}(\overline\theta)<0$.
On a neighborhood of that zero, $m_{0,h}'(\theta)\leq-2j_h<0$.
Moreover,
\begin{equation}
\inf_{\theta\in\Theta}A_h(\theta)\geq a_h>0.
\label{eq:local-cdf-overlap}
\end{equation}
Any numerical CDF clipping threshold is less than $a_h/2$.
\item The split is independent of the observations, and
$|I_k|/n\to\pi_k\in(0,1)$ for each fold. Training, model selection, and
synthetic sampling for a fold use no observations from its holdout set.
\item The nuisance is either equation~\eqref{eq:empirical-local-nuisance}
or the calibrated learner in equation~\eqref{eq:local-calibration} with a
fixed finite pseudo-count (50 in the implementation). Synthetic maxima are
finite almost surely; their number may be fixed. The same weight is used for
$A_h$ and $G_h$ and does not depend on $\theta$.
\item The variance
$V_h=\operatorname{Var}\{\psi_h(B;\theta_{0,h},\eta_{0,h})\}$ is positive.
The numerical routine refines its bracket to score residual
$\varepsilon_n=o(n^{-1/2})$ and uses the total derivative of
$\widehat\Psi_n$; its rejection threshold is at most $j_h$.
\end{enumerate}
\end{assumption}

These conditions state a single-crossing restriction on a specified search
interval, not an assumption that an arbitrary root-selection routine is
consistent. They imply that $\theta_{0,h}$ uniquely maximizes $R_h$ on
$\Theta$. Positive variance and overlap concern the maximum distribution at
this target, and need not hold uniformly over a growing number of bidders.

\begin{proposition}[Primitive uniform rates]
\label{prop:empirical-local-rate}
Under Assumption~\ref{ass:fixed-h}, put $\rho_n=\sqrt{\log n/n}$. For the
empirical-local learner, and also for its fixed-pseudo-count calibration,
\begin{equation}
\max_{k\leq K}\sup_{\theta\in\Theta}
\left\{\|\widehat\eta_h^{(-k)}-\eta_{0,h}\|
+\|\partial_\theta\widehat\eta_h^{(-k)}
-\partial_\theta\eta_{0,h}\|\right\}=O_p(\rho_n).
\label{eq:empirical-local-rate}
\end{equation}
Calibration changes the empirical nuisance and its first derivative by
uniformly $O(n^{-1})$.
Consequently, $\widehat\Psi_n$ and its total derivative converge uniformly to
$m_{0,h}$ and $m_{0,h}'$, respectively.
\end{proposition}

\begin{theorem}[Fixed-bandwidth inference]
\label{thm:fixed-h-dml}
Under Assumption~\ref{ass:fixed-h}, the probability of the stated inference
failures tends to zero. On successful runs the selected root is consistent,
$\widehat J_h=\widehat\Psi_n'(\widehat\theta_h)\to_pJ_h$, and
\begin{equation}
\sqrt n(\widehat\theta_h-\theta_{0,h})
=\frac{1}{\sqrt n}\sum_{i=1}^n\varphi_{i,h}+o_p(1),
\qquad
\varphi_{i,h}=-J_h^{-1}
\psi_h(B_i;\theta_{0,h},\eta_{0,h}).
\label{eq:fixed-h-linearization}
\end{equation}
The empirical variance of the estimated influence values is consistent for
$V_h/J_h^2$. The usual studentized statistic is asymptotically standard normal,
and its Wald interval has asymptotic coverage $1-\gamma$ for
$\theta_{0,h}$, not for $r_0$.
\end{theorem}

The total derivative includes the dependence of both nuisance functions on
$\theta$. If $T_h=(S_h,K_h)^{\mathsf T}$ and a dot denotes this derivative,
\begin{equation}
\frac{d\psi_h}{d\theta}
=m_\theta+\nabla_\eta m_h^{\mathsf T}\dot T_h
+\{\nabla_\eta m_\theta+\nabla_\eta^2m_h\dot\eta\}^{\mathsf T}(T_h-\eta).
\label{eq:total-score-derivative}
\end{equation}
Here $\dot S_h=K_h$ and
$\dot K_h=-(\theta-B)K_h/h^2$. A centered finite difference is also consistent
if its step tends to zero, remains inside $\Theta$, and is applied to the same
smooth pooled score; the uniform derivative result supplies this conclusion.
Fixed numerical tolerances in archived pilots are finite-sample
approximations, not asymptotic guarantees. Failure frequencies are reported
alongside coverage rather than silently treating rejected cases as estimates.

\begin{proposition}[Full-sample plug-in equivalence]
\label{prop:plugin-equivalence}
Let $\widetilde\eta_h(\theta)=n^{-1}\sum_{i=1}^nT_h(B_i;\theta)$ and let
$\widetilde\theta_h$ be the unique downward root of
$m_h(\theta,\widetilde\eta_h(\theta))$ on $\Theta$, refined to the same
residual tolerance. Under Assumption~\ref{ass:fixed-h}, it has the influence
function in equation~\eqref{eq:fixed-h-linearization}, and
\begin{equation}
\sqrt n(\widehat\theta_h-\widetilde\theta_h)\to_p0.
\end{equation}
Thus empirical-local DML, fixed-pseudo-count calibrated DML, and the
same-kernel full-sample plug-in have the same first-order asymptotic variance.
\end{proposition}

This equivalence does not rule out finite-sample differences. It does rule out
attributing a first-order variance improvement to the WGAN component under
these assumptions. The proof uses the empirical part of the calibrated
learner and makes no consistency claim for the standalone WGAN. For a pure
WGAN score, local function and derivative consistency, a negligible remainder
in equation~\eqref{eq:anisotropic-remainder}, and control of synthetic-sample
error require separate verification. A low training critic objective does not
supply those conditions. At fixed $h$, a rate for the true Wasserstein distance
can control the Lipschitz kernel moments; their constants deteriorate as $h$
shrinks, and a fitted critic value is not such a population rate.

\subsection{Confidence sets and conservative envelopes}
\label{sec:score-confidence-sets}

The Wald interval in Theorem~\ref{thm:fixed-h-dml} uses the score slope and
variance evaluated at the estimated root. A second construction retains the
nonlinear score curve and evaluates its variance at each candidate reserve.
For $i\in I_k$, write
\[
\widehat\psi_{i,h}(\theta)
=\psi_h\!\left(B_i;\theta,\widehat\eta_h^{(-k)}(\theta)\right),
\qquad
\widehat V_n(\theta)
=\frac{1}{n-1}\sum_{i=1}^n
\{\widehat\psi_{i,h}(\theta)-\widehat\Psi_n(\theta)\}^2.
\]
This is the sample variance of the pooled observation-level scores, not the
variance of the $K$ fold means. Define the studentized score and its continuous
acceptance set by
\begin{equation}
T_n(\theta)=\frac{\sqrt n\,\widehat\Psi_n(\theta)}
{\sqrt{\widehat V_n(\theta)}},
\qquad
\mathcal C_{n,1-\gamma}
=\{\theta\in\Theta: |T_n(\theta)|\leq z_{1-\gamma/2}\}.
\label{eq:score-confidence-set}
\end{equation}
The set is defined at candidates where the statistic is well defined. On a
failed point-estimation run no confidence set is reported; for coverage
accounting the set is taken to be empty. The critical value
is fixed before the experiment; it is 1.96 in the 95 percent comparisons below.
Neither the point estimate nor the nuisance-learning procedure is changed.

\begin{proposition}[Fixed-bandwidth score-set coverage]
\label{prop:score-set-coverage}
Under Assumption~\ref{ass:fixed-h},
$T_n(\theta_{0,h})\Rightarrow\mathcal N(0,1)$, and the continuous set in
equation~\eqref{eq:score-confidence-set} satisfies
\[
\Pr\{\theta_{0,h}\in\mathcal C_{n,1-\gamma}\}\longrightarrow1-\gamma.
\]
If the set is nonempty, define its conservative interval envelope as
\begin{equation}
\mathcal E_{n,1-\gamma}
=[\inf\mathcal C_{n,1-\gamma},\sup\mathcal C_{n,1-\gamma}],
\label{eq:score-envelope}
\end{equation}
and set $\mathcal E_{n,1-\gamma}=\varnothing$ when
$\mathcal C_{n,1-\gamma}=\varnothing$. Then
$\mathcal C_{n,1-\gamma}\subseteq\mathcal E_{n,1-\gamma}$ for every sample,
and
\[
\liminf_{n\to\infty}
\Pr\{\theta_{0,h}\in\mathcal E_{n,1-\gamma}\}\geq1-\gamma.
\]
These claims concern the fixed-bandwidth target and the mathematical sets,
not exact finite-sample coverage or a certificate for a finite search grid.
\end{proposition}

The proof in Appendix~\ref{app:score-set-proof} uses the same orthogonal
remainder and variance consistency as the Wald result. Inverting pointwise
tests does not require a simultaneous confidence band over $\Theta$: coverage
is the nonrejection event at the one true parameter. The argument does not
establish robustness to weak identification, growing bidder counts, or
shrinking bandwidths.

The full acceptance set can be disconnected even when the estimated score
has a unique zero. Candidate-dependent studentization may produce a distant
accepted region where the estimated variance is large. I retain every
detected component, rather than selecting the component containing
$\widehat\theta_h$. Such selection could remove covered parameter values.
The interval envelope fills the rejected gaps between components. It is
therefore reported separately and is not called the original score-test
confidence set or an improvement in estimation precision.

The implementation evaluates 281 equally spaced candidates on $\Theta$,
adds the estimated root, and refines detected acceptance-boundary crossings
by Brent's method with tolerance $10^{-11}$. It records the components without
joining gaps. Separate flags identify more than one component and contact
with either search boundary, using tolerance $10^{-9}$ only for the flag.
The boundaries are not moved by this tolerance. A boundary flag refers to the
prespecified search interval, not to a data-identified support endpoint.
Nonfinite candidate moments or nonpositive score variance are recorded as
inference failures; they are not silently dropped from the search. In
contrast, disconnection and boundary contact are reporting flags, not
computational failures. Failure rates use all replications, while reported
lengths condition on a successfully constructed set. For a disconnected set,
its length is the sum of component lengths; the envelope length includes the
intervening gaps. Finer grids check numerical sensitivity but cannot certify
the absence of every arbitrarily narrow component.

\subsection{Diagnostic extensions}

To diagnose a failure of the WGAN-based nuisance, I also evaluate two versions
of the identical score. The empirical-local version replaces generated maxima
with observed maxima from the corresponding training fold when estimating
$(A_h,G_h)$. The oracle version supplies their population values from the known
Monte Carlo DGP. Agreement between these two versions, together with failure of
the WGAN version, isolates the problem in nuisance approximation rather than in
equation~\eqref{eq:orthogonal-reserve-score}.

Inference for the exact reserve adds a separate bandwidth problem. The
fixed-bandwidth theorem does not supply a triangular-array result. Under local
smoothness, Gaussian regularization produces a leading approximation bias of
order $h^2$, while the density component of the score has stochastic scale of
order $(nh)^{-1/2}$. A heuristic undersmoothing requirement is therefore
$nh^5\to0$, together with $nh\to\infty$. For a power rule $h_n\propto
n^{-\alpha}$, this suggests $1/5<\alpha<1$. These rate inequalities do not
guarantee stable finite-sample root selection, especially because the maximum-
bid CDF is small near the reserve. The simulation below reports both coverage
of $\theta_{0,h_n}$ and coverage of $r_0$ rather than treating them as the same
target.

To remove the leading smoothing bias without forcing $h_n$ into the unstable
region, suppose the selected population root admits the local expansion

\begin{equation}
\theta_{0,h}=r_0+b_2h^2+O(h^4).
\end{equation}

For a fixed ratio $\lambda>1$, define the Richardson estimator

\begin{equation}
\widehat r_{\mathrm R}
=\frac{\lambda^2\widehat\theta_h-
\widehat\theta_{\lambda h}}{\lambda^2-1}.
\label{eq:richardson-reserve}
\end{equation}

The $O(h^2)$ term cancels in its population analogue. Both estimates use the
same folds and observations, so their dependence must enter inference. If
$\widehat\varphi_{i,h}$ and $\widehat\varphi_{i,\lambda h}$ are their estimated
influence values, the corrected influence value is

\begin{equation}
\widehat\varphi_{i,\mathrm R}
=\frac{\lambda^2\widehat\varphi_{i,h}-
\widehat\varphi_{i,\lambda h}}{\lambda^2-1},
\qquad
\widehat{\mathrm{se}}(\widehat r_{\mathrm R})
=\frac{\widehat{\mathrm{sd}}(\widehat\varphi_{i,\mathrm R})}{\sqrt n}.
\end{equation}

The implementation sets $\lambda=\sqrt{2}$, so
$\widehat r_{\mathrm R}=2\widehat\theta_h-
\widehat\theta_{\sqrt{2}h}$. It selects the smoother root first and uses it as
the reference for the smaller-bandwidth root. This continuation rule reduces
arbitrary root switching but does not eliminate multiple-root uncertainty.

The following extension is conditional: the joint linear representation,
higher-order bias expansion, and covariance conditions are additional
assumptions, not consequences of Theorem~\ref{thm:fixed-h-dml}. In particular,
the differentiability imposed on $F_0$ in Section~3 is not enough by itself to
justify the fourth-order expansion.

\begin{proposition}[Conditional exact-reserve Richardson inference]
\label{prop:richardson-exact}
Let $h_n\to0$ and fix $\lambda>1$. Suppose the roots at $h_n$ and
$\lambda h_n$ have a joint linear representation with remainder
$o_p(s_{\mathrm R,n}/\sqrt n)$ after the Richardson combination and a valid
triangular-array central limit theorem. Suppose their covariance estimator is
consistent, the relevant population roots are locally unique and separated,
and
\begin{equation}
\theta_{0,h}=r_0+b_2h^2+b_4h^4+o(h^4).
\end{equation}
If $s_{\mathrm R,n}^2$ is the variance of the Richardson influence value and
\begin{equation}
\frac{\sqrt n h_n^4}{s_{\mathrm R,n}}\to0,
\label{eq:richardson-rate}
\end{equation}
then the covariance-aware Richardson statistic, centered at $r_0$ and divided
by $s_{\mathrm R,n}/\sqrt n$, converges to a standard normal distribution.
When $s_{\mathrm R,n}\asymp h_n^{-1/2}$, a sufficient bandwidth condition is
$nh_n\to\infty$ and $nh_n^9\to0$.
\end{proposition}

Proposition~\ref{prop:richardson-exact} isolates what the finite-sample
experiment must check: local nuisance rates along the bandwidth path, root
separation, covariance across bandwidths, and residual bias relative to the
standard error. Proofs are in the appendix.

The remaining $O(h^4)$ term can be diagnosed with a third bandwidth. Strengthen
the local expansion to
\begin{equation}
\theta_{0,h}=r_0+b_2h^2+b_4h^4+O(h^6)
\end{equation}
and let $\widehat r_{\mathrm R}(h)$ denote
equation~\eqref{eq:richardson-reserve}. The adjacent extrapolation is
$\widehat r_{\mathrm R}(\lambda h)$, constructed from $\lambda h$ and
$\lambda^2h$. Since
\begin{equation}
r_{\mathrm R}(\lambda h)-r_{\mathrm R}(h)
=(\lambda^4-1)\{r_{\mathrm R}(h)-r_0\}+O(h^6),
\end{equation}
the absolute residual-bias diagnostic is
\begin{equation}
\widehat B_h=
\frac{|\widehat D_h|+1.96\,\widehat{\mathrm{se}}(\widehat D_h)}
{\lambda^4-1},
\qquad
\widehat D_h=\widehat r_{\mathrm R}(\lambda h)
-\widehat r_{\mathrm R}(h).
\label{eq:residual-bias-bound}
\end{equation}
The standard error of $\widehat D_h$ again combines observation-level
influence values across all three bandwidths and therefore retains their
covariance. I report the bias-envelope interval
\begin{equation}
\left[\widehat r_{\mathrm R}(h)-1.96\widehat{\mathrm{se}}_{\mathrm R}
-\widehat B_h,\;
\widehat r_{\mathrm R}(h)+1.96\widehat{\mathrm{se}}_{\mathrm R}
+\widehat B_h\right].
\label{eq:bias-envelope}
\end{equation}

For root robustness, I enumerate every grid crossing at $h$, connect each to
the nearest root at $\lambda h$ and then $\lambda^2h$, construct
equation~\eqref{eq:bias-envelope} for every continuation chain, and take the
union of those intervals. This retains disconnected components rather than
silently selecting one root. The procedure is a finite-sample diagnostic, not
an honest uniform confidence set: its interpretation requires a stable even-
order bias expansion, adequate grid coverage, and reliable local Gaussian
approximations along each root chain.
